%%
%% This is file `sample-acmcp.tex',
%% generated with the docstrip utility.
%%
%% The original source files were:
%%
%% samples.dtx  (with options: `all,journal,acmcp')
%% 
%% IMPORTANT NOTICE:
%% 
%% For the copyright see the source file.
%% 
%% Any modified versions of this file must be renamed
%% with new filenames distinct from sample-acmcp.tex.
%% 
%% For distribution of the original source see the terms
%% for copying and modification in the file samples.dtx.
%% 
%% This generated file may be distributed as long as the
%% original source files, as listed above, are part of the
%% same distribution. (The sources need not necessarily be
%% in the same archive or directory.)
%%
%%
%% Commands for TeXCount
%TC:macro \cite [option:text,text]
%TC:macro \citep [option:text,text]
%TC:macro \citet [option:text,text]
%TC:envir table 0 1
%TC:envir table* 0 1
%TC:envir tabular [ignore] word
%TC:envir displaymath 0 word
%TC:envir math 0 word
%TC:envir comment 0 0
%%
%% The first command in your LaTeX source must be the \documentclass
%% command.
%%
%% For submission and review of your manuscript please change the
%% command to \documentclass[manuscript, screen, review]{acmart}.
%%
%% When submitting camera ready or to TAPS, please change the command
%% to \documentclass[sigconf]{acmart} or whichever template is required
%% for your publication.
%%
%%
\documentclass[sigconf,review,nonacm]{acmart}
%%
%% \BibTeX command to typeset BibTeX logo in the docs
\AtBeginDocument{%
  \providecommand\BibTeX{{%
    Bib\TeX}}}

%% Rights management information.
\setcopyright{none}
\copyrightyear{2026}
\acmYear{2026}
\acmDOI{XXXXXXX.XXXXXXX}
\acmConference[WSDM '26]{Proceedings of the 19th ACM International Conference on Web Search and Data Mining (WSDM)}{February 22--26, 2026}{Boise, ID}

% acmart already loads hyperref; avoid reloading.
\usepackage{fontawesome5}
\usepackage{float} % for [H] exact placement

\begin{document}

\title{LLM-to-Map: Conversational Co-Simulation of Power and Transportation Networks via LLM-Driven Map Control}

\author{Marouane Benbrahim}
\email{mb4194@msstate.edu}
\affiliation{%
  \institution{Mississippi State University}
  \state{Mississippi}
  \country{USA}
}

\author{Zonghan Zhang}
\email{zz239@msstate.edu}
\affiliation{%
  \institution{Mississippi State University}
  \state{Mississippi}
  \country{USA}
}

\author{Zhiqian Chen}
\email{zchen@cse.msstate.edu}
\affiliation{%
  \institution{Mississippi State University}
  \state{Mississippi}
  \country{USA}
}

\renewcommand{\shortauthors}{Benbrahim et al.}
\acmArticleType{Review}
% \keywords{Traffic-Electricity Co-Simulation, LLM}

\begin{abstract}
We present \textbf{LLM-to-Map}, a conversational interface for controlling and exploring a co-simulation of urban infrastructures. The system integrates SUMO for microscopic traffic modeling and PyPSA for electrical grid simulation, synchronized through an event-driven coordination engine that captures cascading effects between domains. A progressive time system enables continuous simulation flow (08:00 $\rightarrow$ 08:01...) with intelligent vehicle quota management maintaining realistic traffic density. An LLM agent powered by GPT-4 and LangChain provides natural language interaction, translating user queries into simulator actions and visualization commands. LLM-to-Map supports fault-tolerant geospatial search, infrastructure control (e.g., triggering substation failures, activating vehicle-to-grid responses), and interactive map navigation. The Mapbox visualization layer renders disruptions and recovery processes in real-time. We showcase three scenarios: (1) a grid outage cascading into traffic disruptions, (2) stress on substations from high EV charging demand, and (3) V2G-enabled resilience. By combining adaptive co-simulation with conversational control, LLM-to-Map lowers the barrier to studying cross-network resilience. Check our demo at \href{https://github.com/XGraph-Team/Map-LLM}{\faGithub\ Github.com/XGraph-Team/Map-LLM}.
\end{abstract}

\maketitle

\section{Introduction}
Urban infrastructures rarely operate in isolation. Transportation networks depend on electricity for traffic signals and EV charging, while power grids rely on mobility patterns and road access for operation and recovery. Disruptions in one domain can cascade into another: substation failures can disable intersections, and traffic congestion can delay grid repair. Such interdependencies have been recognized in resilience research~\cite{ouyang2014review, buldyrev2010catastrophic}, yet tools for exploring them interactively remain scarce.
Domain-specific simulators like SUMO~\cite{lopez2018sumo} and PyPSA~\cite{brown2018pypsa} excel within their domains, and prior co-simulation frameworks such as GEMINI~\cite{panossian2023architecture} and V2Sim~\cite{qian2024v2sim} have shown promise coupling both networks, particularly under growing EV influence. However, these platforms often require technical expertise and restrict interaction to fixed parameters or scripted runs.

Meanwhile, advances in large language models (LLMs) open the door to natural, conversational interfaces for complex systems. LLM agents have been explored in smart city decision support~\cite{kalyuzhnaya2025llm}, human–AI collaboration in urban planning~\cite{han2025large}, and real-time UI integration~\cite{pozdniakov2024large}, with surveys highlighting their broader potential in multi-agent control~\cite{guo2024large}. These works suggest that natural language interaction can make sophisticated systems accessible to both experts and non-experts.
We introduce \textbf{LLM-to-Map}, a system that combines co-simulation of power and traffic networks with an LLM-driven conversational interface. The core architecture integrates SUMO for microscopic traffic simulation and PyPSA for electrical distribution modeling, synchronized by a coordination engine that propagates events (e.g., grid outages, EV charging loads) across domains. On top of this simulation layer, an LLM agent interprets natural language instructions, mapping them into simulator actions and providing real-time visualization updates. The Mapbox interface provides intuitive rendering of disruptions, recoveries, and cascading behaviors.

Users can issue commands such as "turn off the Penn Station substation", "zoom to Times Square", or "activate vehicle-to-grid support" and observe impacts on both grid and traffic systems, lowering barriers for non-experts while exposing details for resilience research.
The demo highlights three illustrative scenarios:
(1) \textbf{Grid-induced traffic disruption}, where substation failures propagate into mobility bottlenecks; (2) \textbf{EV charging stress}, where surges in demand threaten grid stability; and (3) \textbf{V2G stabilization}, where EVs discharge energy to restore grid balance.

Together, these scenarios illustrate how AI-assisted co-simulation can reveal cross-domain feedback loops and resilience strategies in real time. Beyond transportation and power, the modular design supports extension to communication, supply chains, or other infrastructures, positioning the framework as a generalizable platform for future research.
By coupling infrastructure simulation with AI-driven interaction, our demo offers a new paradigm for exploring, reasoning about, and acting upon interconnected urban systems. It enables natural, human-centered experimentation with cascading failures and resilience measures, advancing intelligent systems that bridge data, interaction, and decision-making.

\section{System Overview}
Prior co-simulation platforms mainly target single-domain or transmission/distribution grid studies~\cite{mihal2022smart}, or focus on EV-grid interactions in isolation~\cite{balogun2023ev}. Our system differs by coupling microscopic traffic simulation with detailed grid modeling, emphasizing cross-domain cascading behaviors.
The \textbf{LLM-to-Map} framework consists of three core layers: (1) \textbf{Domain simulators}, (2) \textbf{Coordination engine}, and (3) \textbf{Interaction and visualization layer}.

\noindent\textbf{Domain Simulators}
We integrate two open-source simulators: SUMO for microscopic transport modeling~\cite{lopez2018sumo} and PyPSA for electrical network analysis~\cite{brown2018pypsa}. SUMO models a central Manhattan road network (34th–59th Street), including hundreds of edges, intersections, and traffic signals with actuated logic. Electric vehicles are explicitly represented, with battery state-of-charge tracked to trigger charging events at mapped stations. PyPSA represents eight major distribution substations (600--1000 MVA) connected to transmission nodes, distribution transformers, and downstream loads. The model accounts for both static baseline demand and dynamic EV charging loads, solving power flow every five seconds to capture operational impacts.

\noindent\textbf{Event-Driven Coordination with Progressive Time.}
A Python coordination engine synchronizes SUMO and PyPSA through a shared clock. At each step, events are propagated across domains: substation failures disable traffic lights and force vehicles into cautious flow, EV charging sessions aggregate into additional substation loads, and line overloads trigger alerts or shedding. This event-driven design preserves causality and supports cascading effects in real time. The architecture is modular, enabling extension to other infrastructures such as communication or water systems.

\noindent\textbf{LLM-to-Map Interaction Layer.}
The novel contribution of this demo lies in the \textbf{LLM-to-Map layer}, which connects simulation backends with user-facing interaction. A GPT-4~\cite{achiam2023gpt} + LangChain agent processes natural language instructions and maps them to infrastructure or visualization commands. Example interactions include:
\begin{itemize}
    \item "Turn off the Penn Station substation" $\rightarrow$ grid node failure cascades through PyPSA and SUMO,
    \item "Show me Times Square" $\rightarrow$ camera pans and zooms in Mapbox,
    \item "Activate vehicle-to-grid rescue" $\rightarrow$ EVs discharge energy into the nearest substation.
\end{itemize}
Intents are categorized into \emph{location search}, \emph{infrastructure control}, \emph{data query}, and \emph{navigation}.

\begin{table}[H]
\centering
\caption{Geospatial Search Capabilities}
\small
\begin{tabular}{lccc}
\toprule
\textbf{Feature} & \textbf{LLM-to-Map} & \textbf{Google Maps} & \textbf{ArcGIS} \\
\midrule
Typo Tolerance & Yes ($\leq$2 edits) & Limited & No \\
Context Tracking & Multi-turn & No & No \\
Live Data Mining & Yes & Partial & Manual \\
Natural Language & Full LLM & Keywords & No \\
Infrastructure Data & Power+Traffic & General & Custom \\
Learning Curve & Conversational & Low & High \\
\bottomrule
\end{tabular}
\end{table}

\noindent\textbf{Visualization and Playback.}
The Mapbox interface renders the Manhattan testbed in 2D with grid nodes, traffic signals, and EV stations. Disruptions are animated with flashing outages, EV discharges, and cinematic camera transitions.
\vspace{-0.3cm}
\begin{figure}[H]
    \centering
    \includegraphics[width=0.95\linewidth]{co-sim.pdf}
    \caption{Co-simulation structure}
    \label{fig:structure}
\end{figure}
\vspace{-0.3cm}
\section{Demonstration Scenarios}
We illustrate the system through three application-driven scenarios highlighting power-transportation interdependence. Scenarios are launched via natural language commands with effects visualized in real time.
\par\noindent
\begin{minipage}{\columnwidth}
\centering
\includegraphics[width=0.9\columnwidth]{Screenshot.png}
\captionof{figure}{System Overview Interface}
\label{fig:interface}
\end{minipage}
\par\noindent\textbf{Scenario 1: Grid-Induced Charging Disruptions.}
In the first scenario, a distribution substation is taken offline using a chatbot command (e.g., "turn off Penn Station substation"). Because EV charging stations are mapped to their nearest substations, the outage immediately disables all charging ports in that service area. Vehicles in SUMO attempting to charge at those stations are forced to reroute to the nearest operational station. This rerouting creates localized queues, with circling behavior until a port becomes available. Vehicles that cannot secure a charging slot in time risk battery depletion and may become stranded on the network.

PyPSA reflects the loss of load from the failed substation and the redistribution of charging demand to neighboring substations. Users can observe the cascading effects in real-time as the grid adapts to the outage and vehicles scramble for alternative charging options. To restore functionality during the simulation, users can type "restore Penn Station" to bring the specific substation back online, or use "restore all" to simultaneously restore all failed substations across the network. These restoration commands instantly re-enable the associated charging infrastructure, allowing queued vehicles to redistribute and resume normal charging patterns.

This scenario highlights how a single grid outage can cascade into widespread charging disruptions, producing ripple effects in both grid demand and mobility patterns. The ability to dynamically fail and restore substations provides users with an interactive way to explore infrastructure resilience and understand the critical interdependencies between power grid reliability and electric vehicle operations.
\vspace{-0.3cm}
\begin{figure}[H]
\centering
\includegraphics[width=0.75\columnwidth]{scenario1.png}
\caption{Scenario 1: Grid-induced charging disruption. (1) Pre-command state: Penn Station substation online, normal grid operations; (2) Post-command impact: Penn Station substation offline, area affected, power disrupted.}
\label{fig:scenario1}
\end{figure}
\vspace{-0.3cm}
\noindent\textbf{Scenario 2: EV Charging Stress.}
The second scenario explores how high electric vehicle penetration stresses distribution substations, with a particular focus on cascading failures during extreme events. Vehicles in SUMO spawn with battery states between 25--30\%, immediately triggering charging needs since EVs seek charging when their battery level drops to 30\% or below. By default, each vehicle searches for the closest operational EV station to its current location. Each charging station includes a mix of DC fast chargers and Level 2 ports; when demand exceeds capacity, vehicles are forced to circle or risk stranding if state-of-charge falls below critical levels.

To demonstrate critical infrastructure vulnerability, users can trigger a blackout scenario by typing "Blackout scenario" in the chatbot interface. This command simulates a widespread grid failure where all substations fail except one, creating an extreme stress test of the system. When activated, the camera automatically zooms out to display the entire map, providing a bird's-eye view as substations fail sequentially, their indicators turning from green to red one by one across the network. After the cascade completes, the camera smoothly zooms in and focuses on the sole remaining operational substation, then continuously pans around the area to capture multiple perspectives of the unfolding crisis. The progressive time system maintains vehicle quotas throughout, continuously spawning replacements as vehicles complete trips, sustaining pressure on the single operational station.

With only one station left, all EVs converge on the same location. The EV station rapidly reaches its maximum capacity as vehicles arrive seeking urgent charging. The rolling camera captures the desperate pattern of remaining vehicles circling the area, waiting for an available charging spot. This creates a visible queue of vehicles continuously looping around the station, with the camera following their circular paths and documenting vehicles becoming stranded as their batteries deplete completely before securing a charging position. The coordination engine aggregates active charging sessions into substation loads, producing voltage drops and thermal overloads in PyPSA. Users can issue queries such as "show charging near Times Square" to zoom into stations and observe the growing queues, circling patterns, and stranded vehicles in real-time. This scenario emphasizes the feedback loop from mobility behavior into grid stress, illustrating potential risks in future electrified transportation systems and highlighting the critical importance of redundancy in charging infrastructure to prevent cascading failures during emergency conditions.
\vspace{-0.3cm}
\begin{figure}[H]
\centering
\includegraphics[width=0.95\columnwidth]{scenario2.png}
\caption{Scenario 2: Blackout stress test. (1) Normal grid operation: all 8 substations functional, distributed EV charging load; (2) Critical infrastructure failure: single substation operational (1/8); (3) System response: all EVs rerouting to single station, capacity reached (20/20 occupied), queue formation visible; (4) Critical impact: multiple vehicles stranded (battery depleted), unable to access overloaded charging station.}
\label{fig:scenario2}
\end{figure}
\vspace{-0.3cm}
\noindent\textbf{Scenario 3: Vehicle-to-Grid Rescue.}
In the third scenario, vehicle-to-grid (V2G) is activated as a resilience measure. Users can test this capability by typing "run V2G scenario" in the chatbot, which triggers a Times Square substation failure that leaves people stuck in an elevator, demonstrating the human impact of grid failures. When the substation fails, the system identifies nearby EVs with battery levels of 70\% or higher, sending them notifications to participate in the V2G protocol. These notifications offer monetary compensation in exchange for discharging energy back into the grid.

For this scenario, vehicles spawn with 70\% or higher battery levels to ensure sufficient participation in the V2G response. The camera automatically zooms in on the failed Times Square substation, providing a detailed view of the crisis and recovery process. Users can observe EVs responding to the V2G request, connecting to discharge points, and collectively injecting power back into the grid. The aggregated power injection raises voltages, restores the substation operation—freeing the trapped elevator occupants—and brings traffic signals back online. Mapbox visualizations show affected EVs rerouting to V2G-enabled stations and discharging energy, with restored signals flashing back to green as power returns.

Once the grid is stabilized, the chatbot provides a comprehensive summary of the entire V2G session, including: total power restored to the grid, number of vehicles that participated, individual and collective monetary compensation earned by EV owners, and the exact power requirements that were met through the distributed response. Beyond simulation, this scenario echoes ongoing utility-level pilots in Texas and California, where EV fleets have been leveraged as distributed backup resources. In our demo, it also acts as a countermeasure to Scenario 1, closing the loop between disruption and recovery while demonstrating how EVs can transform from grid stressors to grid stabilizers during emergencies.
\vspace{-0.3cm}
\begin{figure}[H]
\centering
\includegraphics[width=0.95\columnwidth]{scenario3.png}
\caption{Scenario 3: V2G emergency response. (1) Normal operation: Times Square substation online; (2) Substation failed: 18 people trapped in elevators; (3) V2G active: EVs (>70\% charge) discharging power for payment; (4) V2G success: power restored, elevators and traffic signals online.}
\label{fig:scenario3}
\end{figure}
\vspace{-0.3cm}
\noindent\textbf{Interactive Exploration.}
While each scenario can be played in scripted form, the demo is designed for exploration. Attendees may combine scenarios (e.g., induce an outage during peak EV charging) or issue free-form queries such as ”zoom to Grand Central and show traffic impact." The LLM agent interprets these conversational inputs, executes corresponding simulator actions, and returns synchronized visualizations. This interaction paradigm transforms the demo into a live laboratory where cascading failures and resilience strategies can be studied directly.



\section{Implementation}
The \textbf{LLM-to-Map} system is implemented as a modular Python-based architecture that couples simulation backends, a coordination engine, and a natural language interface.

\noindent\textbf{Backend and Coordination.}
The backend is built with Python Flask and a custom event-driven coordination engine. SUMO and PyPSA run as separate processes, each advancing at its native time step. The coordination engine maintains a shared clock (100 ms base) and propagates events across simulators: grid outages disable mapped traffic lights, EV charging sessions contribute dynamic loads, and vehicle-to-grid discharges inject power into failed substations. This design allows synchronization of heterogeneous simulators in real time.

\noindent\textbf{LLM Agent and Query Processing.}
Natural language interaction is powered by GPT-4 through the OpenAI API, orchestrated with LangChain. User queries are first preprocessed with fuzzy-string matching (Levenshtein distance, $\leq$2 edits) to handle location name variations. The LLM then classifies intents into \emph{location search}, \emph{infrastructure control}, \emph{data query}, or \emph{navigation}, and generates structured API calls to the coordination engine.

\section{Discussion and Outlook}
The \textbf{LLM-to-Map} system demonstrates how conversational interfaces can bridge the gap between technical co-simulation tools and accessible urban analysis. By integrating SUMO, PyPSA, and Mapbox under a common event-driven engine, the system captures bidirectional effects between power and transportation networks, while an LLM agent lowers the barrier for interaction. Through natural language, users can trigger outages, observe cascading failures, and activate vehicle-to-grid responses without requiring domain expertise or manual configuration.

While the current implementation focuses on Manhattan and supports English queries, the architecture is generalizable. Any city with OpenStreetMap and utility data can be incorporated, and multilingual LLMs can expand accessibility to broader audiences. Beyond transportation and power, the modular design allows additional infrastructures, such as communication, water, or supply chains, to be added as plugins. This opens opportunities for studying more complex cross-network interactions in the future.

From a research perspective, the system contributes to both web search and urban computing: it applies fuzzy retrieval and conversational state tracking to geospatial data, and it enables interactive what-if analysis of resilience strategies. For practitioners, it offers a hands-on decision support tool that can make cascading risks and mitigation strategies visible in real time. We envision LLM-to-Map evolving into a collaborative platform where researchers, operators, and citizens can explore smart city resilience together.

\bibliographystyle{ACM-Reference-Format}
\bibliography{z}

\end{document}
\endinput
%%
%% End of file `sample-acmcp.tex'.
